\documentclass[11pt]{article}

\usepackage[T1]{fontenc}
\usepackage[utf8]{inputenc}
\usepackage{lmodern}
\usepackage[margin=1.1in]{geometry}
\usepackage{amsmath,amssymb,amsthm,mathtools}
\usepackage{enumitem}
\usepackage{microtype}
\usepackage{xcolor}
\usepackage{booktabs}
\usepackage{array}
\usepackage{titlesec}
\usepackage{hyperref}

\titleformat{\section}
  {\normalfont\LARGE\bfseries}{\thesection}{1em}{}
\titleformat{\subsection}
  {\normalfont\Large\bfseries}{\thesubsection}{1em}{}

\hypersetup{
  pdftitle={Cubic Collision Geometry and the S3 Sheet},
  pdfauthor={Chloe van der Vlugt},
  colorlinks=true,
  linkcolor=blue!55!black,
  citecolor=green!40!black,
  urlcolor=blue!60!black
}

\newtheorem{theorem}{Theorem}[section]
\newtheorem{proposition}[theorem]{Proposition}
\newtheorem{lemma}[theorem]{Lemma}
\newtheorem{corollary}[theorem]{Corollary}
\theoremstyle{remark}
\newtheorem*{remark}{Remark}
\theoremstyle{definition}
\newtheorem{definition}[theorem]{Definition}

\DeclareMathOperator{\Spec}{Spec}
\DeclareMathOperator{\Frac}{Frac}
\DeclareMathOperator{\Gal}{Gal}
\DeclareMathOperator{\Hom}{Hom}
\DeclareMathOperator{\Aut}{Aut}
\DeclareMathOperator{\Ram}{Ram}
\DeclareMathOperator{\Norm}{Norm}
\DeclareMathOperator{\Ann}{Ann}
\DeclareMathOperator{\core}{core}

\newcommand{\A}{\mathbb A}
\newcommand{\C}{\mathbb C}
\newcommand{\Obs}{\operatorname{Obs}}
\newcommand{\id}{\mathrm{id}}
\newcommand{\qedbox}{\hfill$\square$}
\newcommand{\leanpath}[1]{\texttt{\nolinkurl{#1}}}

\newenvironment{leanfiles}
  {\par\medskip\begingroup\footnotesize\raggedright\noindent
   \textit{Lean correspondence.}\ }
  {\par\endgroup\medskip}

\title{Cubic Collision Geometry and the \(S_3\) Sheet}
\author{Chloe van der Vlugt}
\date{\today}

\begin{document}

\maketitle

\begin{abstract}
For a complex polynomial self-map
\(F\colon\A^n_{\C}\to\A^n_{\C}\), we study the scheme of pairs
\((x,y)\) satisfying \(F(x)=F(y)\).  It contains the diagonal \(x=y\);
write \(C_F\) for its coordinate ring.  The
kernel \(\Obs(F)\) of restriction to the diagonal vanishes exactly when
\(F\) is a polynomial automorphism.
Thus a three-dimensional Keller map is a counterexample to \(JC(3)\)
exactly when its collision scheme does not collapse to the diagonal.  We
show that every possible generic-degree-three counterexample lies in
the nonnormal cubic branch: its residual collision sheet is the normal
closure, whose Galois group is \(S_3\).

Put
\(B=\C[F_1,\dots,F_n]\), \(K=\Frac(B)\), and
\(L=\C(x_1,\dots,x_n)\).  For a dominant map of generic degree three,
\(K\otimes_BC_F\simeq L\otimes_KL\).  The diagonal remains as
the first factor \(L\), alongside a quadratic residual \(L\)-algebra
\(D_\alpha\).  In general, \(D_\alpha\) is split precisely when the
cubic extension \(L/K\) is normal, and connected precisely when it is
nonnormal.  The classical Galois case of the Jacobian conjecture excludes
the split branch for a Keller map: normality would force \(F\) to be an
automorphism and hence force \([L:K]=1\).  Consequently every
generic-degree-three Keller map, if one exists, has, with \(N\) denoting
the normal closure of \(L/K\),
\[
  D_\alpha\simeq_L N,\qquad
  K\otimes_BC_F\simeq L\otimes_KL\simeq L\times N,
  \qquad
  \Gal(N/K)\simeq S_3.
\]
It is therefore a counterexample to \(JC(3)\).  This is a structural
description of the generic-degree-three counterexample locus, not an
existence theorem or a classification of counterexamples of other generic
degrees.  A final dimension-independent normalization result points toward
the boundary problem reserved for a planar sequel.
\end{abstract}

\begin{center}
  \small Licensed under
  \href{https://creativecommons.org/licenses/by/4.0/}
    {CC BY 4.0}.
\end{center}

\tableofcontents

\section{Introduction}
\label{sec:introduction}

Let \(F\colon\A^n_{\C}\to\A^n_{\C}\) be a complex polynomial
self-map.  The basic geometric object of this paper is the scheme of
ordered pairs \((x,y)\) with the same image:
\[
  \operatorname{Coll}(F)
  :=\A^n_{\C}\times_{\A^n_{\C}}\A^n_{\C},
\]
where both maps to the middle copy are \(F\).  It contains the diagonal
\(x=y\), which we retain as a distinguished component.  We ask how the remaining
generic sheets are organized around it.

Our end goal is to describe where counterexamples to the complex
three-dimensional Jacobian conjecture can occur.  For a Keller self-map
of \(\A^3_{\C}\), being a counterexample is equivalent to having a
nonempty residual collision scheme.  In generic degree three, classical
Galois rigidity eliminates the normal cubic branch altogether.  Hence
every possible Keller map in this degree class is necessarily nonnormal,
its quadratic residual sheet is connected, and that sheet is the normal
closure with Galois group \(S_3\).  This is the conclusion toward which
the paper is organized.

For a polynomial map \(F\), the collision and diagonal ideals satisfy
\(I_R\subseteq I_\Delta\).  The quotient
\[
  \Obs(F)=I_\Delta/I_R\subseteq C_F=S/I_R
\]
is the kernel of diagonal evaluation.  Consequently,
\[
  \Obs(F)=0
  \quad\Longleftrightarrow\quad
  I_R=I_\Delta
  \quad\Longleftrightarrow\quad
  F\text{ is a polynomial automorphism}.
\]
This gives a single affine ideal whose support is the failure of the
collision relation to collapse to the diagonal.

The generic fiber makes the field-theoretic content transparent.  Put
\[
  B=\C[F_1,\dots,F_n],\qquad
  K=\Frac(B),\qquad
  L=\C(x_1,\dots,x_n).
\]
Under the usual dominance and generic-finiteness hypotheses, base change
to \(K\) identifies the collision algebra with
\[
  K\otimes_BC_F\simeq L\otimes_KL,
\]
and diagonal evaluation with tensor multiplication.  A primitive element
of a finite separable extension then exhibits the marked root as the first
factor and the residual conjugate algebra as the second.

Suppose now that \([L:K]=3\).  The marked-root decomposition splits
\(L\otimes_KL\) into its diagonal factor and a quadratic residual
\(L\)-algebra \(D_\alpha\).  For a general cubic map, this algebra is
split in the normal branch and connected in the nonnormal branch.  For a
Keller map, the split branch cannot occur: the classical Galois case of
the Jacobian conjecture would make \(F\) an automorphism and force generic
degree one.  Thus \(D_\alpha\) is connected, it is the normal closure
\(N\), and
\[
  L\otimes_KL\simeq L\times N,
  \qquad
  \Gal(N/K)\simeq S_3.
\]
The diagonal factor selects one root, while the residual factor supplies a
second; in a cubic extension those two roots determine the third.  The
faithful action on the three conjugate sheets is what produces \(S_3\).
The nonzero residual factor makes every generic-degree-three Keller map a
counterexample to \(JC(3)\), if such a map exists.  We do not claim that
every possible counterexample has generic degree three.

Section~\ref{sec:general-construction} develops the collision ideal,
fiber-product interpretation, generic base change, and marked-root
decomposition.  Section~\ref{sec:cubic-s3-obstruction} identifies the
two cubic residual branches and shows that a Keller map can occur only in
the nonnormal, \(S_3\)-symmetric branch.  Section
\ref{sec:further-directions} records the dimension-independent
normalization picture that motivates a planar sequel.

\section{General construction}
\label{sec:general-construction}

The organizing question is how the failure of a polynomial self-map to be
one-to-one is encoded by the part of its fiber square lying off the diagonal.
The collision ideal isolates that part scheme-theoretically; generic base
change then identifies its field-level factors with conjugate embeddings.
All constructions in this section are dimension-independent.

\subsection{The two ideals and their obstruction}
\label{sec:collision-ideals}

Throughout the paper, the base field is \(\C\).  Let
\[
  F=(F_1,\dots,F_n)\colon
  X=\A^n_{\C}\longrightarrow\A^n_{\C}
\]
be a polynomial map, and let
\[
  S=\C[x_1,\dots,x_n,y_1,\dots,y_n].
\]
The map \(F\) is Keller if
\[
  \det JF\in\C^\times.
\]
Define
\[
  I_R(F)=\bigl(F_i(x)-F_i(y)\bigr)_{i=1}^{n},
  \qquad
  I_\Delta=(x_i-y_i)_{i=1}^{n},
  \qquad
  \Obs(F):=I_\Delta/I_R(F).
\]
The quotient \(\Obs(F)\) is naturally an ideal of the collision ring
\(C_F:=S/I_R(F)\); we call it the obstruction ideal.
For a fixed map \(F\), we abbreviate \(I_R(F)\) to \(I_R\).

\begin{proposition}[Canonical containment]
\label{prop:collision-contained-in-diagonal}
For every polynomial map \(F\),
\[
  I_R(F)\subseteq I_\Delta.
\]
\end{proposition}

\begin{proof}
For \(H\in\C[x_1,\dots,x_n]\), the difference \(H(x)-H(y)\) belongs to
\((x_1-y_1,\dots,x_n-y_n)\).  This follows by induction on \(H\), or by
telescoping each monomial one coordinate at a time.  Apply the statement
to each coordinate \(F_i\).
\end{proof}

Let \(A:=\C[x_1,\dots,x_n]\).  Diagonal evaluation induces a surjective ring
map
\[
  \bar\mu_F\colon C_F\longrightarrow A,
  \qquad
  [s]_{I_R}\longmapsto s(x,x).
\]

\begin{proposition}[The diagonal kernel]
\label{prop:obstruction-kernel}
As ideals of \(C_F\),
\[
  \ker(\bar\mu_F)
  =
  \Obs(F)
  =
  I_\Delta/I_R.
\]
Consequently,
\[
  \ker(\bar\mu_F)=0
  \quad\Longleftrightarrow\quad
  \Obs(F)=0
  \quad\Longleftrightarrow\quad
  I_R=I_\Delta.
\]
\end{proposition}

\begin{proof}
The kernel of diagonal evaluation \(S\to A\) is \(I_\Delta\).  Passing
through \(S/I_R\) therefore identifies the displayed kernel with
\(I_\Delta/I_R\).  Since \(I_R\subseteq I_\Delta\),
\[
  I_\Delta/I_R=0
  \quad\Longleftrightarrow\quad
  I_R=I_\Delta.
\]
\end{proof}

\subsection{The fiber-product interpretation}
\label{sec:fiber-product}

Let
\[
  A=\C[x_1,\dots,x_n],\qquad
  B=\C[F_1,\dots,F_n]\subseteq A.
\]
Let \(X=\Spec A\), let \(Y=\Spec B\), and let
\[
  f_F\colon X\longrightarrow Y
\]
be the morphism induced by \(B\hookrightarrow A\).
The homomorphisms
\(\iota_x,\iota_y\colon A\rightrightarrows C_F\), defined by
\(\iota_x(a)=[a(x)]\) and \(\iota_y(a)=[a(y)]\), agree on \(B\) because
\([F_i(x)]=[F_i(y)]\) in \(C_F\).  Their common restriction equips
\(C_F\) with a \(B\)-algebra structure.

\begin{theorem}[Fiber-product interpretation]
\label{thm:collision-tensor-product}
There is a canonical \(B\)-algebra equivalence and its dual
affine-scheme equivalence
\[
  C_F\cong A\otimes_B A,
  \qquad
  R_F:=\Spec C_F\cong X\times_YX.
\]
Under the first equivalence, \(\bar\mu_F\) is tensor multiplication
\[
  A\otimes_BA\longrightarrow A,\qquad a\otimes a'\longmapsto aa'.
\]
Its dual morphism is the diagonal
\[
  \Delta_{X/Y}\colon X\longrightarrow X\times_YX.
\]
\end{theorem}

\begin{proof}
A map \(C_F\to T\) to a commutative \(\C\)-algebra \(T\) is equivalent
to a pair of maps \(f,g\colon A\rightrightarrows T\) satisfying
\[
  f(F_i)=g(F_i)\qquad(1\leq i\leq n).
\]
Since the \(F_i\) generate \(B\), this is precisely the universal
property of \(A\otimes_BA\).  Applying \(\Spec\) to this equivalence and
using
\[
  \Spec(A\otimes_BA)
  \cong
  \Spec A\times_{\Spec B}\Spec A
\]
gives the asserted fiber-product equivalence.
The statements about multiplication and the diagonal follow on pure
tensors.
\end{proof}

\[
  R_F(\C)
  =
  \{(u,v)\in X(\C)\times X(\C):F(u)=F(v)\}.
\]
Inside
\[
  X\times X=\Spec S,
\]
the ideals \(I_R\) and \(I_\Delta\) define, respectively,
\(X\times_YX\) and the image of \(\Delta_{X/Y}\).  The
contravariance of \(I\mapsto V(I)\) therefore gives the equivalences
\[
  I_R\subseteq I_\Delta
  \quad\Longleftrightarrow\quad
  \Delta_{X/Y}(X)\subseteq X\times_YX
\]
and
\[
  I_R=I_\Delta
  \quad\Longleftrightarrow\quad
  X\times_YX=\Delta_{X/Y}(X)
\]
as closed subschemes of \(X\times X\).  Thus the algebraic and
geometric gaps are two descriptions of the same scheme-theoretic
containment.

\subsection{The residual collision scheme}
\label{sec:residual-collision-scheme}

For every polynomial map \(F\), define
\[
  I_{\mathrm{off}}:=I_R:I_\Delta,
  \qquad
  C_F^\circ:=S/I_{\mathrm{off}},
  \qquad
  R_F^\circ:=\Spec C_F^\circ.
\]
This definition does not require the diagonal to be clopen.

\begin{corollary}[Off-diagonal emptiness criterion]
\label{cor:off-diagonal-empty-criterion}
For every polynomial map \(F\), the following are equivalent:
\[
  R_F^\circ=\varnothing
  \quad\Longleftrightarrow\quad
  C_F^\circ=0
  \quad\Longleftrightarrow\quad
  \Obs(F)=0
  \quad\Longleftrightarrow\quad
  I_R=I_\Delta.
\]
\end{corollary}

\begin{proof}
Affine emptiness gives \(R_F^\circ=\varnothing\) if and only if
\(C_F^\circ=0\), and the latter holds if and only if
\(I_R:I_\Delta=S\).  This colon is the unit ideal exactly when
\(I_\Delta\subseteq I_R\); together with the canonical containment
\(I_R\subseteq I_\Delta\), this is equivalent to equality.
Proposition~\ref{prop:obstruction-kernel} identifies that equality with
\(\Obs(F)=0\).
\end{proof}

\subsection{Generic fibers and marked-root decomposition}
\label{sec:generic-galois-bridge}

Let \(F\) be dominant and generically finite, and let
\[
  A=\C[x_1,\dots,x_n],
  \qquad
  B=\C[F_1,\dots,F_n]\subseteq A.
\]
Let
\[
  K=\Frac(B),
  \qquad
  L=\Frac(A).
\]
Let
\[
  \theta\colon K\otimes_BA\longrightarrow L,
  \qquad
  \lambda\otimes a\longmapsto\lambda a.
\]

\begin{theorem}[Generic collision bridge]
\label{thm:generic-collision-base-change}
The map \(\theta\) is an isomorphism.  Consequently, there is a
\(K\)-algebra isomorphism
\[
  K\otimes_BC_F
  \cong
  L\otimes_KL,
\]
and the following square commutes:
\[
  \begin{array}{ccc}
    K\otimes_BC_F
      &\xrightarrow{\ 1\otimes\bar\mu_F\ }&
    K\otimes_BA\\
    {\scriptstyle\sim}\downarrow
      &&\downarrow{\scriptstyle\theta\ \sim}\\
    L\otimes_KL
      &\xrightarrow{\ \mu_L\ }&
    L.
  \end{array}
\]
\end{theorem}

\begin{proof}
Let \(U=B\setminus\{0\}\).  Since \(K=U^{-1}B\), localization gives
\[
  K\otimes_BA
  \xrightarrow{\sim}
  U^{-1}A,
  \qquad
  \frac{b}{u}\otimes a\longmapsto\frac{ba}{u}.
\]
Since \(A=\C[x_1,\dots,x_n]\) is a domain and \(B\subseteq A\),
\[
  B\setminus\{0\}\subseteq A\setminus\{0\}.
\]
Hence
\[
  U^{-1}A\hookrightarrow\Frac(A)=L,
  \qquad
  \frac{a}{u}\longmapsto\frac{a}{u}.
\]
Composing this injection with the displayed isomorphism gives
\(\theta\), so \(\theta\) is injective.

Since \(F\) is generically finite, \(L/K\) is a finite extension.  The
image of \(U^{-1}A\) in \(L\) is
\[
  K[x_1,\dots,x_n].
\]
Since each \(x_i\) is algebraic over \(K\), this \(K\)-domain is
finite-dimensional over \(K\), hence is a field.  Since it contains
\(A\), it contains \(\Frac(A)=L\); the reverse inclusion is already
given by its embedding into \(L\).  Therefore
\[
  U^{-1}A=L,
\]
and \(\theta\) is an isomorphism.

By Theorem~\ref{thm:collision-tensor-product}, identify
\(C_F\) with \(A\otimes_BA\).  There is a \(K\)-algebra isomorphism
\[
  \beta\colon
  K\otimes_B(A\otimes_BA)
  \xrightarrow{\sim}
  (K\otimes_BA)\otimes_K(K\otimes_BA)
\]
defined on pure tensors by
\[
  \beta\bigl(\lambda\otimes(a\otimes a')\bigr)
  =
  (\lambda\otimes a)\otimes(1\otimes a').
\]
Its inverse is
\[
  (\lambda\otimes a)\otimes(\lambda'\otimes a')
  \longmapsto
  \lambda\lambda'\otimes(a\otimes a').
\]
The balancing relations over \(B\) and \(K\) show that these formulas
are well-defined, and the two displayed maps are mutually inverse.
Composing \(\beta\) with \(\theta\otimes_K\theta\) gives
\[
  \Phi_F\colon K\otimes_BC_F\xrightarrow{\sim}L\otimes_KL.
\]

It remains to identify the diagonal map.  Under
\(C_F\cong A\otimes_BA\), the homomorphism \(\bar\mu_F\) is
multiplication.  Therefore, on a pure tensor,
\[
  \begin{aligned}
  (\mu_L\circ\Phi_F)
    \bigl(\lambda\otimes(a\otimes a')\bigr)
    &=(\lambda a)a'
     =\lambda aa',\\
  \bigl(\theta\circ(1\otimes\bar\mu_F)\bigr)
    \bigl(\lambda\otimes(a\otimes a')\bigr)
    &=\theta(\lambda\otimes aa')
     =\lambda aa'.
  \end{aligned}
\]
Pure tensors generate \(K\otimes_BC_F\), so the two homomorphisms are
equal.  This proves the commutative square.
\end{proof}

Suppose now that \(L/K\) is separable and has a power basis
\(L=K(\alpha)\).  Let \(m_\alpha(T)\in K[T]\) be the minimal polynomial
of \(\alpha\), and let \(h_\alpha(T)\in L[T]\) be determined by
\[
  m_\alpha(T)=(T-\alpha)h_\alpha(T)
  \qquad\text{in }L[T].
\]

\begin{theorem}[Marked-root tensor decomposition]
\label{thm:marked-root-tensor-decomposition}
Give \(L\otimes_KL\) its \(L\)-algebra structure through the left tensor
factor.  There is an \(L\)-algebra equivalence
\[
  L\otimes_KL
  \cong
  L\times L[T]/(h_\alpha),
\]
under which tensor multiplication \(L\otimes_KL\to L\) is projection
onto the first factor.
\end{theorem}

\begin{proof}
The power-basis presentation
\[
  L\cong K[T]/(m_\alpha)
\]
and base change from \(K\) to \(L\) give an \(L\)-algebra isomorphism
\[
  L\otimes_KL
  \xrightarrow{\sim}
  L[T]/(m_\alpha),
  \qquad
  1\otimes\alpha\longmapsto T.
\]
Since \(m_\alpha\) is separable,
\(m_\alpha'(\alpha)\neq0\).  Differentiating
\[
  m_\alpha(T)=(T-\alpha)h_\alpha(T)
\]
and evaluating at \(T=\alpha\) gives
\[
  h_\alpha(\alpha)=m_\alpha'(\alpha)\neq0.
\]
Thus \(T-\alpha\) does not divide \(h_\alpha\), so these two polynomials
are coprime in \(L[T]\).  The Chinese remainder theorem now gives
\[
  L[T]/(m_\alpha)
  \xrightarrow{\sim}
  L[T]/(T-\alpha)\times L[T]/(h_\alpha)
  \xrightarrow{\sim}
  L\times L[T]/(h_\alpha).
\]
Under the first isomorphism above, a pure tensor
\(\ell\otimes p(\alpha)\), with \(p\in K[T]\), corresponds to the class
of \(\ell p(T)\).  Tensor multiplication sends it to
\(\ell p(\alpha)\), which is evaluation at \(T=\alpha\).  Under the
Chinese remainder isomorphism, this evaluation map is projection onto
the factor \(L[T]/(T-\alpha)\cong L\).  Hence tensor multiplication is
the first projection.
\end{proof}

Together with Theorem~\ref{thm:generic-collision-base-change}, this gives
\[
  K\otimes_BC_F
  \cong
  L\times L[T]/(h_\alpha),
\]
where the first projection is the generic diagonal.

\subsection{Automorphism detection}
\label{sec:automorphism-detection}

\begin{theorem}[Automorphism criterion]
\label{thm:automorphism-criterion}
For an arbitrary polynomial self-map
\(F\colon\A^n_{\C}\to\A^n_{\C}\), the following are equivalent:
\begin{enumerate}[label=\textup{(\arabic*)}]
  \item \(F\) is a polynomial automorphism;
  \item \(F\) is injective on \(\A^n(\C)\);
  \item \(I_R(F)=I_\Delta\);
  \item \(\Obs(F)=0\).
\end{enumerate}
\end{theorem}

\begin{proof}
Suppose first that \(F\) has polynomial inverse
\(G=(G_1,\dots,G_n)\).  Then
\[
  x_i-y_i
  =
  G_i(F(x))-G_i(F(y))
  \in I_R
\]
for every \(i\), so \(I_\Delta\subseteq I_R\).  Proposition
\ref{prop:collision-contained-in-diagonal} gives equality.  Thus
\textup{(1)} implies \textup{(3)}.

Suppose \(I_R=I_\Delta\), and let
\(\operatorname{ev}_{(a,b)}\colon S\to\C\) be evaluation at
\((a,b)\).  If \(F(a)=F(b)\), then
\[
  I_R\subseteq\ker(\operatorname{ev}_{(a,b)}).
\]
Therefore \(x_i-y_i\in I_\Delta=I_R\) gives \(a_i-b_i=0\) for every
\(i\).  Thus \(a=b\), proving
\textup{(3)}\(\Rightarrow\)\textup{(2)}.  The
Ax--Grothendieck theorem \cite{Ax1969} gives
\textup{(2)}\(\Rightarrow\)\textup{(1)}.  Finally,
Proposition~\ref{prop:obstruction-kernel} gives
\textup{(3)}\(\Longleftrightarrow\)\textup{(4)}.
\end{proof}

\begin{corollary}[The Jacobian conjecture as obstruction vanishing]
\label{cor:jacobian-conjecture-as-vanishing}
For every \(n\geq1\), the \(n\)-dimensional Jacobian conjecture is
equivalent to each of the statements
\[
  \forall F\colon\A^n_{\C}\to\A^n_{\C},
  \qquad
  \det JF\in\C^\times
  \quad\Longrightarrow\quad
  \Obs(F)=0,
\]
and
\[
  \forall F\colon\A^n_{\C}\to\A^n_{\C},
  \qquad
  \det JF\in\C^\times
  \quad\Longrightarrow\quad
  I_R(F)=I_\Delta.
\]
\end{corollary}

\begin{proof}
Apply Theorem~\ref{thm:automorphism-criterion} to every Keller map.
Its equivalent conditions \textup{(1)}, \textup{(3)}, and \textup{(4)}
give the two displayed formulations.
\end{proof}

We will also use the classical Galois case of the conjecture.  It is the
step that rules out the normal cubic branch under the Keller hypothesis.

\begin{theorem}[Classical Galois rigidity]
\label{thm:classical-galois-rigidity}
Let
\[
  F\colon\A^n_{\C}\longrightarrow\A^n_{\C},
  \qquad \det JF\in\C^\times,
\]
and put
\[
  K=\C(F_1,\dots,F_n),
  \qquad
  L=\C(x_1,\dots,x_n).
\]
If \(L/K\) is normal, then \(F\) is a polynomial automorphism and
\(L=K\).
\end{theorem}

\begin{proof}
The Keller condition makes \(L/K\) finite and separable.  Thus normality
makes it Galois, and the conclusion is the classical Galois case of the
Jacobian conjecture due to Campbell, with algebraic treatments by Razar
and Wright \cite{Campbell1973,Razar1979,Wright1981}.
\end{proof}

\subsection{Generic degree one}

\begin{corollary}[Generic degree-one descent]
\label{cor:generic-degree-one-descent}
If
\[
  L=K
  \qquad\text{and}\qquad
  A\text{ is flat over }B,
\]
then
\[
  \Obs(F)=0.
\]
\end{corollary}

\begin{proof}
When \(L=K\), Theorem~\ref{thm:generic-collision-base-change}
identifies the base-changed diagonal with tensor multiplication:
\[
  \ker(1\otimes\bar\mu_F)
  \cong
  \ker\!\left(
    K\otimes_KK\xrightarrow{\mu_K}K
  \right)
  =
  0.
\]
The fiber-product identification of
Theorem~\ref{thm:collision-tensor-product} is
\[
  C_F\xrightarrow{\sim}A\otimes_BA.
\]
Since \(A\) is flat over \(B\), the tensor product \(A\otimes_BA\), and
therefore \(C_F\), is flat over \(B\).  The localization map
\[
  \iota_C\colon C_F\hookrightarrow K\otimes_BC_F
\]
is therefore injective, and
\[
  \iota_C\bigl(\ker\bar\mu_F\bigr)
  \subseteq
  \ker(1\otimes\bar\mu_F)
  =
  0.
\]
Thus \(\ker\bar\mu_F=0\).  Proposition
\ref{prop:obstruction-kernel} identifies
\[
  \ker(\bar\mu_F)=\Obs(F)=I_\Delta/I_R,
\]
so \(\Obs(F)=0\).
\end{proof}

\begin{leanfiles}
\leanpath{CollisionIdeals/General/Collision/CollisionDiagonal.lean},
\leanpath{CollisionIdeals/General/ResidualCollision/Scheme.lean},
\leanpath{CollisionIdeals/General/GenericFiber/CollisionDiagonal.lean}, and
\leanpath{CollisionIdeals/General/Automorphism/Equivalences.lean}.  The Lean form of
Theorem~\ref{thm:generic-collision-base-change} takes surjectivity of
\(\theta\) as a parameter; the proof above derives it from generic
finiteness.
\end{leanfiles}

% The planar specialization is retained in main.tex and reserved for Paper II.

\section{\texorpdfstring{The cubic \(S_3\) obstruction in dimension three}
  {The cubic S3 obstruction in dimension three}}
\label{sec:cubic-s3-obstruction}

We now specialize the dimension-independent constructions of
Section~\ref{sec:general-construction} to generic degree three.  After the
diagonal is exhibited as the first factor, the complementary
generic collision algebra has degree two over \(L\).  For an
arbitrary cubic map it is either split,
corresponding to a normal cubic extension, or connected, corresponding to
a nonnormal cubic extension.  Classical Galois rigidity rules out the
split branch for a Keller map of generic degree three.  The only possible
Keller branch is therefore connected: its residual factor is the normal closure,
it carries \(S_3\)-symmetry, and its survival away from the diagonal is
precisely the failure of \(JC(3)\).

Let \(S_3=\operatorname{Perm}(\{1,2,3\})\) be the symmetric group on
three letters.

\subsection{The cubic extension and its normal closure}
\label{sec:cubic-extension}

Let
\[
  F=(F_1,F_2,F_3)\colon\A^3_{\C}\longrightarrow\A^3_{\C}
\]
be dominant and generically finite.  Let
\[
  A=\C[x_1,x_2,x_3],
  \qquad
  B=\C[F_1,F_2,F_3]\subseteq A,
  \qquad
  K=\Frac(B),
  \qquad
  L=\Frac(A),
\]
and let
\[
  \theta\colon K\otimes_BA\longrightarrow L,
  \qquad
  \lambda\otimes a\longmapsto\lambda a,
\]
be the isomorphism of
Theorem~\ref{thm:generic-collision-base-change}.  Assume that \(F\) has
generic degree three:
\[
  [L:K]=3.
\]
Because \(K\) has characteristic zero, \(L/K\) is separable.  Choose a
primitive element \(\alpha\), so that \(L=K(\alpha)\) with power basis
\(1,\alpha,\alpha^2\).
If \(m_\alpha\in K[T]\) is the minimal polynomial of \(\alpha\), let
\(h_\alpha\in L[T]\) be determined by
\[
  m_\alpha(T)=(T-\alpha)h_\alpha(T)
  \qquad\text{in }L[T].
\]
The marked-root decomposition identifies the algebra remaining after the
diagonal factor with
\[
  D_\alpha:=L[T]/(h_\alpha),
\]
a finite \(L\)-algebra of degree two.
Let \(N/K\) be the normal closure of \(L/K\), containing \(L\) through
the marked inclusion.  Let
\[
  G=\Gal(N/K),
  \qquad
  H=\Gal(N/L).
\]
Restriction of automorphisms gives a bijection
\[
  G/H\xrightarrow{\sim}\Hom_K(L,N),
  \qquad gH\longmapsto g|_L.
\]
Because the conjugates of \(L\) generate its normal closure \(N\), an
element of \(G\) fixing every conjugate embedding fixes \(N\).  Hence the
coset action is faithful and
\[
  \core_G(H)=\bigcap_{g\in G}gHg^{-1}=1.
\]

The residual collision algebra detects which cubic branch occurs.  We
include the short argument in order to fix its marked identification with
the normal closure.

\begin{proposition}[Connected cubic residual factor]
\label{prop:cubic-tensor-decomposition}
The following conditions are equivalent:
\begin{enumerate}[label=\textup{(\roman*)}]
  \item \(\Spec D_\alpha\) is connected;
  \item \(D_\alpha\) is a field;
  \item \(h_\alpha\) is irreducible over \(L\);
  \item \(L/K\) is nonnormal.
\end{enumerate}
When they hold, there is an \(L\)-algebra isomorphism
\[
  \varphi\colon D_\alpha\xrightarrow{\sim}_L N.
\]
Consequently,
\[
  \Psi\colon L\otimes_KL\xrightarrow{\sim}_L L\times N,
  \qquad
  \operatorname{pr}_1\circ\Psi=\mu_L.
\]
\end{proposition}

\begin{proof}
Write
\[
  m_\alpha(T)=T^3+a_2T^2+a_1T+a_0.
\]
The quadratic \(h_\alpha\) is separable.  If it is irreducible, then
\(D_\alpha\) is a field and hence has connected spectrum.  If it is
reducible, it is a product of two distinct linear factors; the Chinese
remainder theorem then gives \(D_\alpha\simeq L\times L\), whose spectrum
is disconnected.  This proves the equivalence of \textup{(i)},
\textup{(ii)}, and \textup{(iii)}.

If \(h_\alpha\) is reducible, it has a root \(\beta\in L\).
Separability gives \(\beta\neq\alpha\), and the third root is
\[
  \gamma=-a_2-\alpha-\beta\in L.
\]
Thus \(m_\alpha\) would split over \(L\).  Since
\(L=K(\alpha)\), the extension \(L/K\) would then be the splitting
field of the separable polynomial \(m_\alpha\), hence normal.  Conversely,
if \(L/K\) is normal, then \(m_\alpha\) splits over \(L\), so
\(h_\alpha\) is reducible.  Hence \textup{(iii)} and \textup{(iv)} are
equivalent.

Assume these equivalent conditions, and let \(\beta\) denote the image of
\(T\) in \(D_\alpha\).  This field contains the three roots
\[
  \alpha,\qquad \beta,\qquad -a_2-\alpha-\beta
\]
of \(m_\alpha\), and is generated over \(L\) by \(\beta\).  It is
therefore the splitting field of \(m_\alpha\), which is the normal
closure \(N\) of \(L/K\).  Choose the resulting \(L\)-algebra
isomorphism
\[
  \varphi\colon D_\alpha\xrightarrow{\sim}_L N.
\]

Let
\[
  \Phi_\alpha\colon
  L\otimes_KL
  \xrightarrow{\sim}_L
  L\times D_\alpha
\]
be the marked-root decomposition of
Theorem~\ref{thm:marked-root-tensor-decomposition}.  It satisfies
\[
  \operatorname{pr}_1\circ\Phi_\alpha=\mu_L.
\]
Composing its second factor with \(\varphi\), let
\[
  \Psi=(\id_L\times\varphi)\circ\Phi_\alpha.
\]
Since \(\id_L\times\varphi\) leaves the first factor unchanged,
\(\operatorname{pr}_1\circ\Psi=\mu_L\).
\end{proof}

\begin{corollary}[The cubic Keller residual is connected]
\label{cor:cubic-keller-residual-connected}
If \(F\) is Keller and \([L:K]=3\), then \(L/K\) is nonnormal and
\(\Spec D_\alpha\) is connected.
\end{corollary}

\begin{proof}
If \(L/K\) were normal, Theorem~\ref{thm:classical-galois-rigidity}
would make \(F\) a polynomial automorphism and therefore give \(L=K\),
contrary to \([L:K]=3\).  Proposition
\ref{prop:cubic-tensor-decomposition} now identifies nonnormality with
connectedness of \(\Spec D_\alpha\).
\end{proof}

\begin{theorem}[\(S_3\)-symmetry of the connected residual sheet]
\label{thm:index-three-core-free-s3}
If \(\Spec D_\alpha\) is connected, then
\[
  \Gal(N/K)\cong S_3.
\]
\end{theorem}

\begin{proof}
Proposition~\ref{prop:cubic-tensor-decomposition} identifies \(N\) with
the degree-two field \(D_\alpha\) over \(L\).  Hence
\[
  |H|=[N:L]=2.
\]
The normal-closure identities recalled above give
\[
  \core_G(H)=1,
  \qquad
  [G:H]=[L:K]=3.
\]
Hence the action of \(G\) on the three cosets of \(H\) gives
\[
  \rho\colon G\longrightarrow\operatorname{Perm}(G/H)\cong S_3,
  \qquad
  \ker\rho=\core_G(H)=1.
\]
Core-freeness therefore makes \(\rho\) injective, so
\[
  |G|\leq |S_3|=6.
\]
The index formula gives
\[
  |G|=[G:H]|H|=3\cdot2=6.
\]
Thus \(|G|=6\), and the injective homomorphism \(\rho\) is an
isomorphism \(G\cong S_3\).
\end{proof}

\subsection{The cubic off-diagonal factor}

Retain the cubic data of Section~\ref{sec:cubic-extension} and assume that
\(F\) is Keller.  Corollary
\ref{cor:cubic-keller-residual-connected} shows that \(D_\alpha\) is
connected and \(L/K\) is nonnormal; Proposition
\ref{prop:cubic-tensor-decomposition} then identifies \(D_\alpha\) with
\(N\).  Thus
\[
  B=\C[F_1,F_2,F_3]\subseteq A=\C[x_1,x_2,x_3],
  \qquad
  K=\Frac(B)\subseteq L=\Frac(A)=K(\alpha)\subseteq N,
\]
where \([L:K]=3\) and \(N/K\) is the normal closure.  Let
\[
  \theta\colon K\otimes_BA\xrightarrow{\sim}L,
  \qquad
  \lambda\otimes a\longmapsto\lambda a.
\]

\begin{corollary}[Generic cubic \(S_3\) collision]
\label{cor:cubic-s3-collision}
There is a \(K\)-algebra isomorphism
\[
  \Psi_F\colon
  K\otimes_BC_F\xrightarrow{\sim}_K L\times N,
  \qquad
  \operatorname{pr}_1\circ\Psi_F
  =
  \theta\circ(1\otimes\bar\mu_F).
\]
It satisfies
\[
  \Psi_F\!\left(
    \ker\bigl(\theta\circ(1\otimes\bar\mu_F)\bigr)
  \right)
  =
  0\times N,
  \qquad
  \Gal(N/K)\cong S_3.
\]
If
\[
  C_{F,K}:=K\otimes_BC_F,
  \qquad
  J_{F,K}:=\ker\bigl(\theta\circ(1\otimes\bar\mu_F)\bigr),
\]
then the intrinsic residual quotient of the generic collision algebra is
\[
  C_{F,K}/\Ann_{C_{F,K}}(J_{F,K})\simeq_K N.
\]
Consequently,
\[
  \Obs(F)\neq0,
  \qquad
  I_R\subsetneq I_\Delta,
  \qquad
  R_F^\circ\neq\varnothing,
  \qquad
  F\notin\Aut_{\mathrm{poly}}(\A^3_{\C}).
\]
\end{corollary}

\begin{proof}
Theorem~\ref{thm:generic-collision-base-change} gives a \(K\)-algebra
isomorphism
\[
  \Phi_F\colon K\otimes_BC_F\xrightarrow{\sim}_K L\otimes_KL
\]
such that
\[
  \mu_L\circ\Phi_F
  =
  \theta\circ(1\otimes\bar\mu_F).
\]
Proposition~\ref{prop:cubic-tensor-decomposition} gives
\[
  \Psi\colon L\otimes_KL\xrightarrow{\sim}_L L\times N,
  \qquad
  \operatorname{pr}_1\circ\Psi=\mu_L.
\]
After restricting \(\Psi\) from \(L\) to \(K\), let
\[
  \Psi_F=\Psi\circ\Phi_F.
\]
It follows that
\[
  \operatorname{pr}_1\circ\Psi_F
  =
  \theta\circ(1\otimes\bar\mu_F),
  \qquad
  \Psi_F\!\left(
    \ker\bigl(\theta\circ(1\otimes\bar\mu_F)\bigr)
  \right)
  =
  0\times N\neq0.
\]
The annihilator of \(0\times N\) in \(L\times N\) is \(L\times0\).
Therefore \(\Psi_F\), followed by the second projection, induces
\[
  C_{F,K}/\Ann_{C_{F,K}}(J_{F,K})\xrightarrow{\sim}_K N.
\]

Suppose that the diagonal-evaluation homomorphism
\(\bar\mu_F\colon C_F\to A\) were injective.  Since \(K\) is a
localization of \(B\), it is flat over
\(B\), so \(1\otimes\bar\mu_F\) would also be injective.
Theorem~\ref{thm:generic-collision-base-change} shows that \(\theta\)
is an isomorphism.  Therefore
\(\theta\circ(1\otimes\bar\mu_F)\) would be injective, contradicting
the nonzero kernel \(0\times N\) above.  Hence
\(\ker(\bar\mu_F)\neq0\), and
Proposition~\ref{prop:obstruction-kernel} gives
\(\Obs(F)\neq0\).  Proposition
\ref{prop:collision-contained-in-diagonal} shows that
\(I_R\subseteq I_\Delta\) strict, and
Corollary~\ref{cor:off-diagonal-empty-criterion} gives
\(R_F^\circ\neq\varnothing\).

Finally, Theorem~\ref{thm:index-three-core-free-s3} identifies
\(\Gal(N/K)\) with \(S_3\).  If \(F\) were a polynomial automorphism,
Theorem~\ref{thm:automorphism-criterion} would give
\(\Obs(F)=0\), contrary to \(\Obs(F)\neq0\).  Since \(F\) is Keller, it is
a dimension-three Jacobian-conjecture counterexample.
\end{proof}

\begin{remark}[Ordered roots]
Let
\[
  m_\alpha(T)=T^3+a_2T^2+a_1T+a_0
\]
have roots \(\alpha_1,\alpha_2,\alpha_3\) in \(N\).  For distinct
\(i,j,k\),
\[
  \alpha_k=-a_2-\alpha_i-\alpha_j,
  \qquad
  K(\alpha_i,\alpha_j)
  =
  K(\alpha_1,\alpha_2,\alpha_3)
  =
  N.
\]
Thus the residual factor, which marks two distinct roots, is already
the ordered-root normal closure.
\end{remark}

\begin{leanfiles}
\leanpath{CollisionIdeals/ComplexThree/Cubic/GaloisGroup.lean},
\leanpath{CollisionIdeals/ComplexThree/Cubic/S3Collision.lean},
\leanpath{CollisionIdeals/ComplexThree/OffDiagonal.lean},
\leanpath{CollisionIdeals/ComplexThree/Cubic/Main.lean}, and
\leanpath{CollisionIdeals/ComplexThree/Statements/JacobianConjecture.lean}.  The
generic product-annihilator quotient is formalized in
\leanpath{CollisionIdeals/General/GenericFiber/DiagonalObstruction.lean}.  The
Lean cubic witness takes the residual-field isomorphism and
\(H\neq1\) as explicit inputs.  In the manuscript these follow from
Proposition~\ref{prop:cubic-tensor-decomposition} and Corollary
\ref{cor:cubic-keller-residual-connected}; the use of classical Galois
rigidity in that corollary is not yet formalized.  Hidden inertia is formalized
separately in
\leanpath{CollisionIdeals/ComplexThree/Cubic/HiddenInertia.lean}, so the cubic
core has no normalization or ramification dependency.
\end{leanfiles}

\begin{remark}[Exact relation with classical complex \(JC(3)\)]
\label{rem:exact-relation-jc3}
The traditional complex three-dimensional Jacobian conjecture asserts that
every polynomial map
\[
  F\colon\A^3_{\C}\longrightarrow\A^3_{\C},
  \qquad
  \det JF\in\C^\times,
\]
is a polynomial automorphism; equivalently, every such map satisfies
\(\Obs(F)=0\).  Corollary~\ref{cor:cubic-s3-collision} gives the following
implication:
\[
  \begin{gathered}
    \det JF\in\C^\times,\qquad [L:K]=3\\
    \Longrightarrow\quad
    \Spec D_\alpha\text{ connected},\qquad
    \Gal(N/K)\cong S_3,\\
    \Longrightarrow\quad \Obs(F)\neq0
    \quad\Longrightarrow\quad
    F\text{ is a counterexample to }JC(3).
  \end{gathered}
\]
The nonvanishing of \(\Obs(F)\) already follows from \([L:K]=3\).
The additional content is that Theorem
\ref{thm:classical-galois-rigidity} excludes the normal cubic branch,
after which the residual algebra itself produces the normal closure and
its \(S_3\)-symmetry.  Consequently, every generic-degree-three Keller map
would disprove \(JC(3)\), while \(JC(3)\) implies that no such map exists.
This section does not construct such a map, and it does not assert that
every possible counterexample to \(JC(3)\) has generic degree three.
\end{remark}

\section{Further directions}
\label{sec:further-directions}

The normal closure also carries a dimension-independent boundary picture,
although none of it is needed for the cubic collision theorem.  Let
\(L/K\) be finite separable, let \(N/K\) be its normal closure, and let
\(Y\) be a normal integral model with function field \(K\).  Write
\[
  G=\Gal(N/K),\qquad H=\Gal(N/L),\qquad
  Z=\Norm_N(Y)\longrightarrow \overline X=\Norm_L(Y)\longrightarrow Y.
\]
Suppose that a normal model \(X\) of \(L\) is an open subscheme of
\(\overline X\) over \(Y\), and that \(X\to Y\) is \'etale.

\begin{proposition}[Hidden inertia on a conjugate sheet]
\label{prop:generic-hidden-inertia}
Let \(E\in Z^{(1)}\) be ramified over \(Y\), with divisorial valuation
\(w_E\), restriction \(v=w_E|_K\), and inertia group \(I_E\).  Then some
\(g\in G\) satisfies
\[
  e\bigl((g^{-1}w_E)|_L/v\bigr)
  =[I_E:I_E\cap gHg^{-1}]>1.
\]
For every such \(g\), the center of \((g^{-1}w_E)|_L\) on
\(\overline X\) lies in \(\overline X\setminus X\).
\end{proposition}

\begin{proof}
Because \(N\) is the normal closure of \(L/K\), the action of \(G\) on
\(G/H\) is faithful, so \(\bigcap_{g\in G}gHg^{-1}=1\).  Since
\(I_E\neq1\), it is not contained in every conjugate of \(H\); the usual
inertia formula gives the displayed index for a suitable \(g\).  If the
corresponding center belonged to \(X\), the \'etale morphism \(X\to Y\)
would force its ramification index to be one, a contradiction.
\end{proof}

In the generic-degree-three Keller situation of
Section~\ref{sec:cubic-s3-obstruction}, the cubic extension is necessarily
nonnormal and \(G\cong S_3\).  Thus any
divisorial ramification in the ordered-root normalization is detected on a
conjugate sheet and forced into the deleted boundary.  This statement is
conditional on the existence of such a ramified divisor; the cubic
collision theorem itself neither requires nor asserts it.

The cubic and planar problems use the same diagonal in opposite
ways.  Here generic degree three forces a nonzero residual factor beside
it.  In the planar sequel, the goal is to prove that the off-diagonal
factor vanishes, so that the collision scheme consists only of the
diagonal.

A sequel will develop this normalization picture for planar Keller maps:
conjugate affine sheets, their boundary and fixed-locus ideals, and the
secant--trace comparison.  Its remaining Keller-specific input is a
uniform landing theorem converting a finite secant--frame lattice into
control of arbitrary boundary sections.  Purity and finite-\'etale
rigidity enter only after that boundary statement is established.

\section{Conclusion}
\label{sec:conclusion}

The collision ideal turns injectivity into a scheme-theoretic equality.  The
obstruction \(\Obs(F)\) is precisely the kernel of diagonal evaluation,
and it vanishes exactly when the collision relation equals the diagonal.
After passage to the generic fiber, that diagonal becomes the
multiplication factor of \(L\otimes_KL\), while the remaining factors
record the other embeddings of \(L\).

In generic degree three, the diagonal persists as the first factor beside
a quadratic residual algebra.  The classical Galois case rules out its
split, normal branch for a Keller map; hence the residual algebra is connected, it is
the normal closure \(N\), its Galois group is \(S_3\), and
\[
  K\otimes_BC_F\simeq L\times N.
\]
Its nonzero residual factor gives \(\Obs(F)\neq0\), so every
generic-degree-three Keller map, if one exists, is a counterexample to
\(JC(3)\) with an \(S_3\)-normal-closure sheet.  The theorem does not
assert the existence of such a map or classify counterexamples of other
generic degrees.
The normalization principle
of Section~\ref{sec:further-directions} explains how the same conjugate
sheet geometry leads to boundary questions, which are reserved for the
planar sequel.

\paragraph{Motivating example.}
The explicit three-dimensional polynomial counterexample announced by
Levent Alpöge \cite{Alpoge2026} and independently formalized in
Isabelle/HOL by Ramos, Hulak, and de Queiroz
\cite{AFPJacobianCounterexample} motivated the cubic \(S_3\) comparison.
The formal verification establishes a
constant nonzero Jacobian determinant and distinct rational points with
the same image.  That particular map is not used in the arguments above:
the dimension-three result is the structural implication from the
generic-degree-three Keller condition.

The formalization uses Lean~4 and Mathlib \cite{Lean4,Mathlib}.

\paragraph{Acknowledgment.}
OpenAI Codex (GPT-5.6 Sol) assisted with Lean proof engineering,
mathematical and bibliographic checking, and manuscript editing.  The
author reviewed and takes responsibility for all statements, proofs,
citations, and formalization code.

\bibliographystyle{alpha}
\bibliography{references}

\end{document}
